\begin{abstract}	
\noindent Employees value jobs in the formal and informal employment differently because of a mismatch in social protection, job protection and flexibility. There remains a need of distinct policy recognition to ensure uniform welfare distribution in the economy. The main objective of this paper is to investigate if a wage difference exists between employees in formal and informal employment based on their individual, job based and employer characteristics in the Nepalese labor market using a sample of full time employed workers. We use data of Nepal Labor Force Survey III. We adopt Oaxaca-Blinder decomposition to decompose in the context of mean wages. We use quantile regression  to capture the heterogeneity in workers' responses to observed attributes and calculate the wage gap at various quantiles of the wage distribution. Finally, we present the conditional quantile decomposition techniques developed by Machado and Mata. The results reveal the existence of wage gap between formal and informal employment; education level, gender, job sector and size of business being key determinants of earning gap. Our finding shows the unexplained effects in wage gap decreases with the movement along higher distribution, implying that the job market in Nepal shows sticky floor effects, with larger discrimination at the bottom of the distribution.\par
\vspace{.1in}
\noindent\textbf{Keywords:} Labor participation, Informal Employment, Wage Gap, Selection bias.\\
\noindent\textbf{JEL Codes:} J08, J31, C21
\bigskip
\end{abstract}